\section{Self-Renormalization of lattice matrix elements}\label{sec:StrtoRe}
\subsection{The Strategy}\label{Strategy}



The non-perturbative matrix elements contain the following important contributions: a) the linear divergence, b) a finite term coming from non-perturbative renormalon physics, and c) the residual contribution encoding intrinsic non-perturbative physics. It is part c) that is required to extract the partonic structure information. Thus, it is important to separate out the latter from the former in a systematic way and with high precision. 

Here we develop a self-renormalization method to extract the residual intrinsic non-perturbative physics and renormalization factor from a matrix element itself without using a different matrix element. The extraction of the residual is much harder than the extraction of the linear divergence term because after we take the logarithm of the matrix element (Eq.~(\ref{eq:logM_s})), the linear divergence term is dominant and the residual is very sensitive to the subtraction. We need to properly take into account the fine features of the data, including logarithmic divergences and discretization errors. 

Based on Eq.~(\ref{eq:ZO_RS_higher}), we modify the fitting function Eq.~(\ref{eq:logM_s}) to be, 
\begin{align}\label{eq:logM}
\ln \mathcal{M}(z,a) = \frac{k z}{a \ln[a \Lambda_{\rm QCD}]} 
+ g(z) + f_{1,2}(z)a \nonumber\\
+ \frac{3 C_{F}}{b_0} \ln \bigg[\frac{\ln [1 /(a \Lambda_{\rm QCD})]}{\ln [\mu / \Lambda_{\rm QCD}]}\bigg]+\ln \bigg[1+\frac{d}{\ln (a \Lambda_{\rm QCD})}\bigg],
\end{align}
where the first term is linearly divergent. $g(z)$ is the residual, which contains the non-perturbative $m_{0}$ effect and the intrinsic non-perturbative physics. $f_{1,2}(z)a$ takes into account the discretization errors, which we allow to be different for the data calculated from MILC and RBC ensembles ($f_{1}(z)$ for MILC and $f_{2}(z)$ for RBC). The last two terms come from the resummation of leading and sub-leading logarithmic divergences, which only affect the overall normalization at different lattice spacings.

We use Eq.~(\ref{eq:logM}) to fit the logarithms to the data for the bare matrix elements for each $z$. The $z$ points can be chosen in 
a large range where the lattice discretization error is small and, at the same time, the statistical error is limited. We fit the dependence on $a$, treating $g(z)$, $f_{1}(z)$, $f_{2}(z)$ as free parameters and fixing $\mu$ at 2 GeV, which is the renormalization scale we choose to use. We will discuss $k$ and $\Lambda_{\rm QCD}$ in detail in Sec.~\ref{ReUncer} and $d$ in Sec.~\ref{sec:d}. We obtain the residual $g(z)$ from the fit.

The $g(z)$ obtained this way does not correspond to what is in the perturbative $\overline{\rm MS}$ scheme because it contains a non-perturbative $m_{0}$ effect. 
We need to eliminate this effect through matching the lattice result to the continuum scheme at short distance $z$. Based on Eq.~(\ref{eq:Reoperator}) and Eq.~(\ref{eq:m0}), we use the following equation within a window $a \ll z < 1/\mu$ to extract $m_{0}$,   
\bea\label{eq:gztoMS}
g(z) - \ln[Z_{\overline{\mathrm{MS}}}(z,\mu, \Lambda_{\overline{\mathrm{MS}}})] = m_{0} z,
\eea
where $Z_{\overline{\mathrm{MS}}}$ is the perturbative matrix element in the $\overline{\mathrm{MS}}$ scheme. For the $O_{\gamma_t}(z)$ matrix elements in the pion, nucleon, and off-shell quark state, we take $Z_{\overline{\mathrm{MS}}}$ at one loop~\cite{Izubuchi:2018srq},
\bea\label{eq:ZMS}
Z_{\overline{\mathrm{MS}}}(z, \mu, \Lambda_{\overline{\mathrm{MS}}}) = 1+\frac{\alpha_{s}(\mu,\Lambda_{\overline{\mathrm{MS}}}) C_{F}}{2 \pi}\bigg[\frac{3}{2} \ln \frac{z^{2} \mu^{2}}{4 e^{-2 \gamma_{E}}}+\frac{5}{2}\bigg]\ , \nonumber\\
\eea
from operator product expansion, where we fix $\mu$ = 2 GeV and $\Lambda_{\overline{\mathrm{MS}}}$ = 0.3 GeV. 
The renormalized physical matrix element is then obtained from
\bea\label{eq:gz-m0}
\mathcal{M}(z)_{R} = \exp[g(z) - m_{0}z],
\eea
and is expected to be consistent with $Z_{\overline{\mathrm{MS}}}$ at small $z$.

One could consider higher-order $\overline{\mathrm{MS}}$ matrix elements for matching, including summing over large logarithms. We choose not to do this here as it does not affect the procedure for the subsequent analysis, and therefore does not change the main conclusions of the paper. However,
to obtain physical results at high precisions, one might need to examine this more carefully. 

We thus define a renormalization factor
\begin{align}\label{eq:ZRm0}
Z(z,a)_{R} = \exp\Big[\frac{k z}{a \ln[a \Lambda_{\rm QCD}]} 
+ m_{0}z + f_{1,2}(z)a \nonumber\\
+\frac{3 C_{F}}{b_0} \ln \bigg[\frac{\ln [1 /(a \Lambda_{\rm QCD})]}{\ln [\mu / \Lambda_{\rm QCD}]}\bigg]+\ln [1+\frac{d}{\ln (a \Lambda_{\rm QCD})}]\Big]\ ,
\end{align}
which includes the discretization error as well. 
All the parameters here are either fixed, fitted or fine-tuned through the renormalization procedure from the matrix element $\mathcal{M}(z,a)$ itself. Dividing the bare matrix element $\mathcal{M}(z,a)$ by $Z(z,a)_{R}$, we get
\bea\label{eq:M/Z}
\mathcal{\widetilde{M}}_{R}(z) = \mathcal{M}(z,a)/Z(z,a)_{R}.
\eea
If our procedure eliminates all the divergences and discretization errors, we should expect that $\mathcal{\widetilde{M}}_{R}(z)$ is $a$-independent and the same as $\mathcal{M}(z)_{R}$ (Eq.~(\ref{eq:gz-m0})). The degree to which this is the case can be used as a test of the self-renormalization procedure.


\subsection{Renormalon Uncertainty}\label{ReUncer}
Now let us turn to the parameter $k$ and $\Lambda_{\rm QCD}$ in Eq.~(\ref{eq:logM}). In Eq.~(\ref{eq:logM}), we have only considered the leading order perturbative term for the linear divergence and neglected higher-order ones. These can be included partially by a proper choice of $\Lambda_{\rm QCD}$. 
For example, the coupling constants with different choices of $\Lambda$ are related to each other by perturbation theory~\cite{Lepage:1992xa}, 
\bea\label{eq:alpha_s_series}
\alpha_{s}(Q,\Lambda_{\rm QCD}) \sim \alpha_{s}(Q,\Lambda) + c_{2} \alpha_{s}^2(Q,\Lambda) \nonumber\\ + c_{3} \alpha_{s}^3(Q,\Lambda) + ...
\eea
In principle, $\Lambda_{\rm QCD}$ depends on the specific lattice action, and one needs to perform multi-loop calculations to relate them. In our analysis, we treat $\Lambda_{\rm QCD}$ as a fitting parameter which, therefore, can effectively take into account some higher-order effects. 

\begin{figure}[tbp]
\centering
\includegraphics[width=9cm]{images/fitting_clover_Pion.pdf}
\caption{Using Eq.~(\ref{eq:logM}) to fit the bare $O_{\gamma_t}(z)$ matrix element in the pion state for the clover action without  HYP  smearing. Blue points are interpolated data and colorful curves are fitted curves for each $z$. The fitted curves are broken lines because $f_{1}(z)$ and $f_{2}(z)$ are allowed to be different. $k$, $\Lambda_{\rm QCD}$, and $d$ are fixed at 7.4 GeV$^{-1}$fm$^{-1}$ (1.46 if dimensionless), 0.09 GeV, and $-1$ respectively. The $\chi^{2}/{\rm d.o.f.}$ of the fitting for each $z$ is listed in the plot legend.
}
\label{fig:fitting_clover_Pion}
\end{figure}

\begin{figure}[tbp]
\centering
\includegraphics[width=9cm]{images/chisqmap_clover_Pion.pdf}
\caption{ 
$\chi^{2}$ map with respect to $k$ and $\Lambda_{\rm QCD}$. $\langle  \chi^{2}/{\rm d.o.f.} \rangle_{z}$ is the average of $\chi^{2}/{\rm d.o.f.}$ among the fitting for each $z$. The ``small-$\chi^{2}$ band'' shows that $k$ and $\Lambda_{\rm QCD}$ are strongly correlated and uncertain for the chosen data.
}
\label{fig:chisqmap_clover_Pion}
\end{figure}

Next we want to discuss the coefficient $k$ of the linear divergence. QCD perturbation theory not only predicts the $1/a$ dependence and linear $z$ dependence (which we have tested in Sec.~\ref{sec:test}), but also the value of $k$. Comparing Eq.~(\ref{eq:delta_m}) and Eq.~(\ref{eq:logM}), we get the one-loop perturbative value of $k$:
\bea\label{eq:k}
k~=~\frac{(2 \pi)^2}{3 b_{0}} = 1.46 = 7.4 {\rm GeV^{-1} fm^{-1}}.
\eea
Next, we check whether this value is consistent with our data.

We use Eq.~(\ref{eq:logM}) to fit the bare $O_{\gamma_t}(z)$ matrix element in the pion state for the clover action without  HYP smearing, treating $g(z)$, $f_{1}(z)$, $f_{2}(z)$ as fitted parameters. If we take ($k$, $\Lambda_{\rm QCD}$) to be one set of values, e.g., (7.4 GeV$^{-1}$fm$^{-1}$, 0.09 GeV), we can perform a fitting for each $z$, as seen in Fig.~\ref{fig:fitting_clover_Pion}. There is a $\chi^{2}/{\rm d.o.f.}$ for the fit at each $z$ and we calculate the average of them, denoted as $\langle \chi^{2}/{\rm d.o.f.} \rangle_{z}$.
Varying ($k$, $\Lambda_{\rm QCD}$) generates a $\chi^{2}$ map, as seen in Fig.~\ref{fig:chisqmap_clover_Pion}. 

The sets of ($k$, $\Lambda_{\rm QCD}$) of small $\chi^{2}$ lie in a band, which we call the ``small-$\chi^{2}$ band''. That means that $k$ and $\Lambda_{\rm QCD}$ are strongly correlated, while there is a large uncertainty for $k$ and $\Lambda_{\rm QCD}$ separately.
If we choose some sets of values from the ``small-$\chi^{2}$ band'' to do the fitting, we find the fitted residuals Exp[$g(z)$] are very different, as seen in Fig.~\ref{fig:Egz}. However, after eliminating the non-perturbative $m_{0}$ effect, the renormalized matrix elements Exp[$g(z) - m_{0} z$] are all the same (see Fig.~\ref{fig:Egz_m0}). Therefore, the uncertainty in $k$ and $\Lambda_{\rm QCD}$ will not significantly influence the final physical result.

\begin{figure}[tbp]
\centering
\includegraphics[width=8cm]{images/Egz.pdf}
\caption{
The fitted residual Exp[$g(z)$] for different sets of ($k$, $\Lambda_{\rm QCD}$) along the ``small-$\chi^{2}$ band''.
}
\label{fig:Egz}
\end{figure}

\begin{figure}[tbp]
\centering
\includegraphics[width=8cm]{images/Egz_m0.pdf}
\caption{
The renormalized matrix element Exp[$g(z)-m_{0} z$] for different sets of ($k$, $\Lambda_{\rm QCD}$) along the ``small-$\chi^{2}$ band''.
}
\label{fig:Egz_m0}
\end{figure}

The reason for this uncertainty lies in the behavior of $\frac{k}{a \ln[a \Lambda_{\rm QCD}]}$ (the term related to the linear divergence in Eq.~(\ref{eq:logM})) in the $a$ range of our data. Varying along the "small-$\chi^2$ band'' effectively shifts $\frac{k}{a \ln[a \Lambda_{\rm QCD}]}$ by a constant $C$, see Fig.~\ref{fig:linearf}. The $Cz$ term is absorbed into $g(z)$ automatically during fitting so there is a large uncertainty for $k$ and $\Lambda_{\rm QCD}$. However, when we eliminate the $m_{0}$ effect, the $Cz$ term has been taken into account,  and the final result does not change. 

\begin{figure}[tbp]
\centering
\includegraphics[width=8cm]{images/LinearTermPlot.pdf}
\caption{ 
The linear divergent term in the $a$ range of our data for different sets of ($k$, $\Lambda_{\rm QCD}$) along the ``small-$\chi^{2}$ band''.
}
\label{fig:linearf}
\end{figure}

Another way to interpret this is that 
the perturbative value of $k$ is consistent with the ``small-$\chi^{2}$'' band 
and therefore is consistent with the fit. As a consequence,
we can fix $k$ at the one-loop perturbative value, 7.4 ${\rm GeV^{-1}fm^{-1}}$ (or 1.46 if dimensionless) in the subsequent analysis. Even with this
choice of $k$, $\Lambda_{\rm QCD}$ can still vary within a reasonable range of $\chi^2$ while the physical result is stabilized by the choice of $m_0$. We call this correlation between
$\Lambda_{\rm QCD}$ and $m_0$ the ``phenomenological renormalon uncertainty'', 
in the sense that $\Lambda_{\rm QCD}$ is characteristic for the different 
truncation/resummation schemes for the perturbation series, but that the non-perturbative $m_0$ compensates the effects of the differences~\cite{Ji:1995tm,Beneke:1998ui,Bauer:2011ws,Bali:2013pla}.

\subsection{Tuning Parameter-$d$ to Match the Continuum Scheme}\label{sec:d}

\begin{figure}[tbp]
\centering
\includegraphics[width=8cm]{images/gz-lnZMS_d0.pdf}
\includegraphics[width=8cm]{images/gz-lnZMS_d.pdf}
\caption{
The fitting to extract $m_{0}$ based on Eq.~(\ref{eq:gztoMS}) for the bare $O_{\gamma_t}(z)$ matrix element in the off-shell quark state for the clover action without HYP smearing. Red points are $g(z) - \ln[Z_{\overline{\mathrm{MS}}}]$, where $g(z)$ is extracted through fitting the data for the bare matrix element using Eq.~(\ref{eq:logM}) and $Z_{\overline{\mathrm{MS}}}$ given in Eq.~(\ref{eq:ZMS}). The blue line is the linear fit of the red points (where the fitting function is $m_{0} z + C_0$). The fitted $m_{0}$ is given in the plot legend. $k$ is fixed at 7.4 GeV$^{-1}$fm$^{-1}$ (1.46 if dimensionless) and $\Lambda_{\rm QCD}$ is chosen as the best fit value 0.118 GeV.
In the upper panel, $d$ is zero and $g(z) - \ln[Z_{\overline{\mathrm{MS}}}]$ is in a linear relationship with $z$ but not proportional to $z$. (It does not go through the origin). In the lower panel, we tune $d$ to make sure that $g(z) - \ln[Z_{\overline{\mathrm{MS}}}]$ is proportional to $z$.
}
\label{fig:gz-lnZMS_d}
\end{figure}

In principle, the parameter $d$ in Eq.~(\ref{eq:logM}) is determined by the lattice action. In our data, there are two different dynamical lattice ensembles (MILC and RBC) and two different valence fermion actions (overlap and clover). Here we treat $d$ as a fine-tuning parameter to make sure that $( g(z) - \ln[Z_{\overline{\mathrm{MS}}}(z,\mu, \Lambda_{\overline{\mathrm{MS}}})] )$ in Eq.~(\ref{eq:gztoMS}) is proportional to $z$ within a window $a \ll z < 1/\mu$.
As shown in Fig.~\ref{fig:gz-lnZMS_d}, this is quite effective.
We have also tested that $d$ has little influence on $\chi^2$ since the effect of tuning $d$ shifts $g(z)$ simply by a constant.

\begin{figure}[tbp]
\centering
\includegraphics[width=8cm]{images/nonpert_clover_quark.pdf}
\caption{
Renormalized $O_{\gamma_t}(z)$ matrix element in the off-shell quark state for the clover action without HYP smearing. Blue points are the renormalized matrix element Exp[$g(z)-m_{0}z$]. The black curve denotes $Z_{\overline{\mathrm{MS}}}$ and the red points are their ratio. $k$ is fixed at 7.4 GeV$^{-1}$fm$^{-1}$ (1.46 if dimensionless) and $\Lambda_{\rm QCD}$ is chosen as the best fitted value 0.118 GeV. The parameter $d$ is fine-tuned to be $-1.184$.
}
\label{fig:nonpert_clover_quark}
\end{figure}

Fig.~\ref{fig:nonpert_clover_quark} shows that the renormalized matrix element Exp[$g(z)-m_{0}z$] (blue) for clover quarks is consistent with $Z_{\overline{\mathrm{MS}}}$ (black) at small $z$. However, at large $z$, there is a significant discrepancy between Exp[$g(z)-m_{0}z$] and $Z_{\overline{\mathrm{MS}}}$. That means that there is a significant non-perturbative effect at large $z$ in the popular RI/MOM and ratio renormalization schemes used previously.
To our knowledge, this is the first time that this difference has been studied 
in the literature. This finding supports the usage of the hybrid renormalization procedure proposed recently~\cite{Ji:2020brr}.

\subsection{Self-Renormalization Procedure}\label{sec:reofpro}

The procedure developed in the previous subsections 
forms our self-renormalization strategy. This procedure can overcome
some of the problems encountered in the renormalization
of the linear divergence due to the numerical uncertainties associated
with exponentially-amplified small errors. 
There are three steps in this process:

1. Use Eq.~(\ref{eq:logM}) to fit the data of the bare matrix elements $\mathcal{M}$. For each $z$ (here $z$ can be chosen in a large range and linear interpolation with respect to $z$ might be needed 
for the data in ln$\mathcal{M}(z,a)$ ), fit the dependence on $a$, treating $\Lambda_{\rm QCD}$ as a global parameter, namely the same for all $z$, and $g(z)$, $f_{1}(z)$, $f_{2}(z)$ as fit parameters, with fixing $k$=7.4 GeV$^{-1}$fm$^{-1}$ (or 1.46 if dimensionless), and $\mu$=2 GeV. Fine-tune $d$ to make sure that $g(z) - \ln[Z_{\overline{\mathrm{MS}}}]$ is proportional to $z$ within a window $a \ll z < 1/\mu$. We obtain the residual $g(z)$ through fitting;

2. Use Eq.~(\ref{eq:gztoMS}) to fit the dependence on $z$ (within a window $a \ll  z < 1/\mu$) to extract $m_{0}$. Then calculate the renormalized matrix element Exp[$g(z)-m_{0}z$] (Eq.~(\ref{eq:gz-m0})), which is considered valid for a large range of $z$;

3. Calculate $\mathcal{M}(z,a)/Z(z,a)_{R}$ (Eq.~(\ref{eq:M/Z})) to see if showing any dependence on the lattice constant $a$ and compare with Exp[$g(z)-m_{0}z$]. 

In step 1 and 2, one can use another way to get $d$ and $m_{0}$: Modify Eq.~(\ref{eq:logM}) through replacing $g(z)$ with ($\ln[Z_{\overline{\mathrm{MS}}}(z,\mu, \Lambda_{\overline{\mathrm{MS}}})] + m_{0} z$), and then use the modified fitting function to fit the bare matrix element at small $z$ to extract $d$ and $m_{0}$. This way should be equivalent to the procedure outlined above and we will not follow it here.

In step 3, we need to calculate $Z(z,a)_{R}$ as well as estimate its error. We shall not propagate the fitting error of $\Lambda_{\rm QCD}$ to $Z(z,a)_{R}$ because the effect of slight variation of $\Lambda_{\rm QCD}$ can be compensated by $m_{0}$ and will not influence our result as we have discussed in Sec.~\ref{ReUncer}.

Since the lattice data we use do not provide systematic errors, we will add a dummy systematic error $\delta_{\rm sys}$ to the data during our fitting~\cite{Zhang:2020rsx}
\bea\label{eq:error}
(\sigma_{\ln \mathcal{M}})_{\rm new} = \sqrt{(\sigma_{\ln \mathcal{M}})_{\rm old}^{2} + (\delta_{\rm sys} a \mu)^{2}},
\eea
where we assume that there is an $a\mu$ dependence as well.
This choice is intended to give more weight to the small lattice spacing 
data in the fitting.

\subsection{Comparing Results from Different Actions and Matrix Elements}

In this subsection, we apply the self-renormalization method to the $O_{\gamma_t}(z)$ matrix 
elements in the pion, nucleon, and off-shell quark state with valence clover and overlap actions 
without HYP smearing. We compare the renormalization factors and physics results. 
We find that apart from the case of the RI/MOM matrix element with clover valence fermion, all other renormalization factors are similar. Despite this difference in renormalization factor, the physical
results are independent of fermion actions. 


\begin{figure*}[tbp]
\centering
\panel{a}{width=0.48\textwidth}{images/overlap_quark_chisqline.pdf}
\panel{b}{height=5.5cm,width=0.48\textwidth}{images/overlap_quark_lnMfitting.pdf}
%\includegraphics[width=8cm]{images/overlap_quark_gz.pdf}
\panel{c}{width=0.48\textwidth}{images/overlap_quark_gzmlnZMS.pdf}
\panel{d}{width=7.5cm}{images/overlap_quark_Expgzmm0z.pdf}
\panel{e}{width=0.48\textwidth}{images/overlap_quark_Expgzmm0zdd.pdf}
\panel{f}{width=0.48\textwidth}{images/overlap_quark_MoZ.pdf}
\caption{
Self-renormalizing the $O_{\gamma_t}(z)$ correlator in the off-shell quark state calculated for the overlap action without HYP smearing, using Eq.~(\ref{eq:logM}). (a) shows $\langle  \chi^{2} \rangle_{z} - [\langle  \chi^{2} \rangle_{z}]_{\rm min}$ with respect to $\Lambda_{\rm QCD}$. The best fitted $\Lambda_{\rm QCD}$ with its error can be estimated here. (b), (c) and (d) show the procedure to get the renormalized matrix element Exp[$g(z)-m_{0}z$]. (b) shows the fitting of the bare matrix element to extract $g(z)$. Blue points are interpolated data and colorful curves are fitted curves. (c) shows the fitting to extract $m_{0}$. (d) shows Exp[$g(z)-m_{0}z$] (blue), compared with $Z_{\overline{\mathrm{MS}}}$ (black). (e) is the comparison of Exp[$g(z)-m_{0}z$] for different $\delta_{\rm sys}$. (f) is the ratio of bare matrix element and renormalization factor (connected with solid lines), compared with Exp[$g(z)-m_{0}z$] (connected with dashed lines).
}
\label{fig:Re_overlap_quark}
\end{figure*}


\begin{figure*}[tbp]
\centering
\panel{a}{width=0.48\textwidth}{images/clover_quark_chisqline.pdf}
\panel{b}{height=5.5cm,width=0.48\textwidth}{images/clover_quark_lnMfitting.pdf}
\panel{c}{width=0.48\textwidth}{images/clover_quark_gzmlnZMS.pdf}
\panel{d}{width=7.5cm}{images/clover_quark_Expgzmm0z.pdf}
\panel{e}{width=0.48\textwidth}{images/clover_quark_Expgzmm0zdd.pdf}
\panel{f}{width=0.48\textwidth}{images/clover_quark_MoZ.pdf}
\caption{
Same as Fig.~\ref{fig:Re_overlap_quark}, with the $O_{\gamma_t}(z)$ correlator in the off-shell quark state calculated for the clover action.
}
\label{fig:Re_clover_quark}
\end{figure*}


\begin{figure*}[tbp]
\centering
\panel{a}{width=0.48\textwidth}{images/overlap_Pion_chisqline.pdf}
\panel{b}{height=5.5cm,width=0.48\textwidth}{images/overlap_Pion_lnMfitting.pdf}
\panel{c}{width=0.48\textwidth}{images/overlap_Pion_gzmlnZMS.pdf}
\panel{d}{width=7.5cm}{images/overlap_Pion_Expgzmm0z.pdf}
\panel{e}{width=0.48\textwidth}{images/overlap_Pion_Expgzmm0zdd.pdf}
\panel{f}{width=0.48\textwidth}{images/overlap_Pion_MoZ.pdf}
\caption{Same as Fig. \ref{fig:Re_overlap_quark}, with quasi-LF correlation in zero-momentum pion state calculated with the overlap fermion action.
}
\label{fig:Re_overlap_Pion}
\end{figure*}


\begin{figure*}[tbp]
\centering
\panel{a}{width=0.48\textwidth}{images/clover_Pion_chisqline.pdf}
\panel{b}{height=5.5cm,width=0.48\textwidth}{images/clover_Pion_lnMfitting.pdf}
\panel{c}{width=0.48\textwidth}{images/clover_Pion_gzmlnZMS.pdf}
\panel{d}{width=7.5cm}{images/clover_Pion_Expgzmm0z.pdf}
\panel{e}{width=0.48\textwidth}{images/clover_Pion_Expgzmm0zdd.pdf}
\panel{f}{width=0.48\textwidth}{images/clover_Pion_MoZ.pdf}
\caption{Same as Fig. \ref{fig:Re_overlap_quark}, with quasi-LF correlation in zero-momentum pion state calculated with the clover fermion action.
}
\label{fig:Re_clover_Pion}
\end{figure*}


\begin{figure*}[tbp]
\centering
\panel{a}{width=0.48\textwidth}{images/clover_Nucleon_chisqline.pdf}
\panel{b}{height=5.5cm,width=0.48\textwidth}{images/clover_Nucleon_lnMfitting.pdf}
\panel{c}{width=0.48\textwidth}{images/clover_Nucleon_gzmlnZMS.pdf}
\panel{d}{width=7.5cm}{images/clover_Nucleon_Expgzmm0z.pdf}
\panel{e}{width=0.48\textwidth}{images/clover_Nucleon_Expgzmm0zdd.pdf}
\panel{f}{width=0.48\textwidth}{images/clover_Nucleon_MoZ.pdf}
\caption{Same as Fig. \ref{fig:Re_overlap_quark}, with quasi-LF correlation in zero-momentum nucleon state calculated with the clover fermion action.
}
\label{fig:Re_clover_Nucleon}
\end{figure*}


Figs.~\ref{fig:Re_overlap_quark} through \ref{fig:Re_clover_Nucleon} show the detailed results of applying the
self-renormalization method for different cases. For the RI/MOM factor with overlap and clover fermions, we analyse eight different lattice spacings from two different types of gauge ensembles. For the pion matrix element calculated with the overlap fermion action, we analyse three different lattice spacings from MILC ensembles. In the clover pion case, we use six different lattice spacings from two different types of gauge ensembles. For the nucleon matrix element with clover fermion action, we analyse five different lattice spacings from MILC ensembles. 

Subfigure (a) in each figure is used to estimate the best fitted value of the global parameter $\Lambda_{\rm QCD}$ as well as its error. For each chosen $\Lambda_{\rm QCD}$, we can use Eq.~(\ref{eq:logM}) to fit the bare matrix element. The fitting for each $z$ gives us a separate $\chi^{2}$. We can calculate the averaged $\chi^{2}$ for different $z$, denoted as $\langle  \chi^{2} \rangle_{z}$. Subfigure (a) shows $\langle  \chi^{2} \rangle_{z} - [\langle  \chi^{2} \rangle_{z}]_{\rm min}$ with respect to $\Lambda_{\rm QCD}$. $\langle  \chi^{2} \rangle_{z} - [\langle  \chi^{2} \rangle_{z}]_{\rm min}$=0 gives us the best fitted $\Lambda_{\rm QCD}$ and $\langle  \chi^{2} \rangle_{z} - [\langle  \chi^{2} \rangle_{z}]_{\rm min}$=1 gives us its error. If the data points for different $z$ were independent, we should use the total $\chi^{2}$ for the fits at different $z$ to estimate the error of $\Lambda_{\rm QCD}$. However, since we have done linear interpolations of the original data points, some of the points are correlated to each other. So here we just use $\langle  \chi^{2} \rangle_{z}$ to estimate the error of $\Lambda_{\rm QCD}$. The range of $z$ starts from 0.18 fm here but not 0.06 fm because the $\chi^{2}$ for $z$ = 0.06 or 0.12 fm is much larger than others. $d$ is fixed to be $-1$ since it has little influence on $\chi^{2}$. 

For $\delta_{\rm sys}=0.002$, the best fitted values of $\Lambda_{\rm QCD}$ are 0.1086(17), 0.1350(21), 0.093(10), 0.086(14), 0.0926(61) GeV for the correlations in the overlap quark, clover quark, overlap pion, clover pion, and clover nucleon, respectively. As one can see, they are quite similar except for the correlation in the clover quark state. While the size of the systematic
error has some effects for the correlation in the quark case, it has little influence on the correlation in the physical 
states. 

In subfigure (b) in each figure, after taking $\Lambda_{\rm QCD}$ to be the best fitted value, we use Eq.~(\ref{eq:logM}) to fit the bare matrix elements to extract $g(z)$. The range of $z$ is taken from 0.06 fm to 1.02 fm with 17 different $z$ values for correlations in the overlap quark case, from 0.06 fm to 1.14 fm, with 19 different $z$ values for correlations in the overlap pion and clover pion cases, from 0.06 fm to 0.96 fm, with 16 different $z$ in the clover quark case, from 0.06 fm to 0.9 fm, and  with 15 different $z$ values in the clover nucleon case. We fine-tune $d$ to make sure that $g(z) - \ln[Z_{\overline{\mathrm{MS}}}]$ is proportional to $z$ within the window 0.06 fm $\leq z \leq$ 0.24 fm. %For $\delta_{\rm sys}=0.002$, 
The fine-tuned $d$ values are $-1.29, -1.35, -1.17, -0.92, -0.97$ in the five cases, respectively. They are all close to $-1$. 

Subfigure (c) in each figure shows the fitting to extract $m_{0}$. For $\delta_{\rm sys}=0.002$, the fitted values of $m_{0}$ are 0.2339(81), 0.7667(83), $-$0.181(24), $-0.363(25)$, and $-0.257(20)$ fm$^{-1}$ in the five cases, respectively. They are all about the order of $\Lambda_{\rm QCD} \sim 0.1$ GeV (about 0.5 fm$^{-1}$), which supports the argument in Sec.~\ref{sec:theory} that $m_{0}$ originates from the non-perturbative renormalon effect. 

In subfigure (d) in each figure, we show our result for the renormalized matrix element Exp[$g(z)-m_{0}z$] in blue points connected by lines. As $z$ increases, Exp[$g(z)-m_{0}z$] increases first, following the perturbative result, 
and then decreases from about $z$ = 0.25 fm on. In all the cases, Exp[$g(z)-m_{0}z$] is consistent with $Z_{\overline{\mathrm{MS}}}$ at small $z$. However, at large $z$, there is a significant discrepancy between Exp[$g(z)-m_{0}z$] and $Z_{\overline{\mathrm{MS}}}$. This means that there is a large non-perturbative effect at large $z$ in the popular RI/MOM and ratio renormalization scheme used previously, which supports the hybrid renormalization procedure proposed recently~\cite{Ji:2020brr}. We have already 
briefly discussed this in Sec.~\ref{sec:d}. 

Subfigure (e) shows the renormalized matrix element Exp[$g(z)-m_{0}z$] for different $\delta_{\rm sys}$. A change in $\delta_{\rm sys}$ does not influence the central values of Exp[$g(z)-m_{0}z$] for most cases except for the clover pion. 
In the clover pion case, the increase of $\delta_{\rm sys}$ can slightly decrease the renormalized matrix element Exp[$g(z)-m_{0}z$]. 

Finally, subfigure (f) shows the ratio of bare matrix element and renormalization factor $\mathcal{M}(z,a)/Z(z,a)_{R}$, compared with the renormalized matrix element Exp[$g(z)-m_{0}z$]. $Z(z,a)_{R}$ here is not extracted from a different matrix elements but from $\mathcal{M}(z,a)$ itself. This subfigure is a consistency check of our method. In all cases, $\mathcal{M}(z,a)/Z(z,a)_{R}$ has little dependence on $a$ within error and is consistent with Exp[$g(z)-m_{0}z$]. 
Therefore, our self-renormalization method can isolate the residual intrinsic physics from the linear divergence, 
renormalon uncertainty, log divergence and discretization error with a reasonable precision. 

\begin{figure}[tbp]
\centering
\includegraphics[width=8.5cm]{images/overlap_clover_Pion_d0_Expgzmm0z.pdf}
\includegraphics[width=8.5cm]{images/overlap_clover_Pion_d0.002_Expgzmm0z.pdf}
\caption{
Comparison of renormalized $O_{\gamma_t}(z)$ correlator in the zero-momentum pion state between overlap and clover 
actions without HYP smearing. Points connected by solid lines are renormalized matrix elements Exp[$g(z)-m_{0}z$]. Black curve is $Z_{\overline{\mathrm{MS}}}$. In the upper panel, $\delta_{\rm sys}=0$. In the lower panel, $\delta_{\rm sys}=0.002$.
}
\label{fig:Re_overlap_clover_Pion}
\end{figure}

Although the ratio $\mathcal{M}(z,a)/Z(z,a)_{R}$ has a good behavior for most of the lattice spacings, it may have large errors for small lattice spacings like 0.0318 fm or 0.0574 fm. Therefore we do not recommend using $\mathcal{M}(z,a)/Z(z,a)_{R}$ as the renormalized matrix element. Since Exp[$g(z)-m_{0}z$] is achieved through our fitting, the
large-error data points have little influence on it. We take this as 
the renormalized matrix element, while using $\mathcal{M}(z,a)/Z(z,a)_{R}$ only 
to test consistency.

We show in Fig.~\ref{fig:Re_overlap_clover_Pion} a comparison of renormalized $O_{\gamma_t}(z)$ correlation in the pion state between overlap and clover actions. Before we add the systematic error, both results show appreciable difference, which
may indicate that two different valence fermions will lead to different results. After we add a small dummy systematic error during the fitting, the correlations in both cases 
lead to similar results. That means that the systematic error is important, and the residual intrinsic non-perturbative physics is independent of valence quark formulations if it is properly included.

\begin{figure}[tbp]
\centering
\includegraphics[width=8.5cm]{images/All_d0.002_Expgzmm0z.pdf}
\caption{
Comparison of renormalized $O_{\gamma_t}(z)$ correlations for all the cases without HYP smearing. Points connected by solid lines are renormalized matrix elements Exp[$g(z)-m_{0}z$]. Black curve is $Z_{\overline{\mathrm{MS}}}$.
}
\label{fig:Re_all}
\end{figure}

Finally, for curiosity, we show in Fig.~\ref{fig:Re_all} the renormalized correlations as well as parameters related to the divergence for all cases we studied. As mentioned before, the fitted $\Lambda_{\rm QCD}$ for overlap quark, overlap pion, clover pion and clover nucleon are similar to each other. But for the clover quark, it is markedly different. We do not yet understand the reason, but it could be due to the combination of chiral symmetry breaking and the linearly-divergent quark mass. We plan to investigate this in the future, but for the time being, at least we understand why we cannot eliminate the linear divergence using RI/MOM for the clover action~\cite{Zhang:2020rsx}. Quite surprisingly, the residual intrinsic non-perturbative 
correlations for overlap quark, pion and nucleon are very similar to each other,
which we also do not fully understand.
%However, this may indicate the correlations in the 
%physical states are independent of valence fermion formulations, a very important conclusion
%if true.
